\section{Local imaginary-time conjugation}

\begin{theorem}[Short imaginary-time conjugation on a chain]
    \label{thm:chain_local_imaginary_time}
    \lean{TNLean.MPS.AreaLaw.norm_exp_conjugation_nearestNeighbour_shift_le_two}
    Let $L,q,c\in\N$, and let $I$ be a set of nearest-neighbour bonds of the
    chain $\{0,\ldots,L-1\}$. For each $i\in I$, let $h_i$ act on
    $\{i,i+1\}$ and satisfy $0\le h_i\le\Id$. Set $H=\sum_{i\in I}h_i$.
    Let $O$ act on $[c,c+2)\cap[0,L)$. For every $z\in\C$ and
    $0\le t\le1/8$,
    \begin{align}
        \bigl\lVert e^{t(H-z\Id)}Oe^{-t(H-z\Id)}\bigr\rVert
        \le 2\lVert O\rVert.
    \end{align}
    This is a nearest-neighbour consequence of the commutator argument
    in~\cite[lines 795--880]{Arad2014Connecting}.
\end{theorem}
\begin{proof}
    Write $O_m=\operatorname{ad}_H^m(O)$. Disjoint supports commute, so $O_m$
    acts on $[\max(0,c-m),c+2+m)\cap[0,L)$. At most $2m+3$ original bonds
    meet this interval. Since $0\le h_i\le\Id$, centering at $\Id/2$ gives
    $\lVert[h_i,X]\rVert\le\lVert X\rVert$. Consequently
    $\lVert O_{m+1}\rVert\le4(m+1)\lVert O_m\rVert$, and induction gives
    $\lVert O_m\rVert\le4^m m!\lVert O\rVert$. Scalar shifts leave every
    commutator unchanged. The Hadamard expansion therefore yields
    \begin{align}
        \bigl\lVert e^{t(H-z\Id)}Oe^{-t(H-z\Id)}\bigr\rVert
         & \le\sum_{m=0}^{\infty}\frac{t^m}{m!}\lVert O_m\rVert
        \le\sum_{m=0}^{\infty}2^{-m}\lVert O\rVert
        =2\lVert O\rVert.
    \end{align}
\end{proof}
